
défini par @@M0@@ si @@M1@@, @@M2@@. Il se factorise par l'inclusion de @@M3@@ dans @@M4@@, ce qui montre que l'homomorphisme de @@M5@@ dans @@M6@@ est injectif, et nous permettra d'identifier le premier à un sous-groupe du second, et de même nous considérerons @@M7@@ comme contenu dans @@M8@@.

On notera que @@M9@@ n'est autre que l'augmentation de @@M10@@.

1.5. On plonge @@M11@@ dans @@M12@@ en envoyant @@M13@@ sur la @@M14@@-ème fonction symétrique des @@M15@@~; en faisant opérer le groupe symétrique @@M16@@ par permutation des @@M17@@, et en utilisant le théorème des polynômes symétriques et (1.3.ii)), on voit que l'homomorphisme de @@M18@@ dans @@M19@@ qu'on en déduit est injectif, ce qui nous permettra encore d'identifier le premier à un sous-groupe du second~:
@@M20@@% source-fix p300 §1.5 (2026-07-03): restored Grothendieck's text vs the paraphrase — removed interpolated «élémentaire», «.»→«;», «qui s'en déduit»→«qu'on en déduit», «permet»→«nous permettra», inlined the H(...) terms, restored dropped intermediate iso H(K[T]^{Sym_n})→~H(K[T])^{Sym_n}.
On note
@@M21@@
l'homomorphisme défini par @@M22@@, et on pose
@@M23@@
Le groupe symétrique permute les morphismes @@M24@@, et laisse donc invariant le morphisme @@M25@@, qui par suite prend ses valeurs dans
